\section{Note on Sources and Delivery}

For an adult parish assembly on the First Sunday of Advent, 29 November 2026, in the 1962 Roman Rite. The homily follows the reviewed studies' interpretation “The Judge whom the faithful long to see.” Jesus tells his disciples in advance whose coming they can expect when the signs begin. The relation between that promise, Paul's encouragement to believers who still quarrel, and the Church's request for rescue from the danger of her own sins is the homily's own. The Gradual passes over the plea for the unjust to be ashamed at the beginning of Ps 24:4 in Challoner and takes up the request for instruction in the same verse. The quarrel and its repair apply that petition and the change of desire Augustine describes at the Communion psalm. The prayers follow the Vatican typical edition of the \work{Missale Romanum}. The Gospel quotation follows Challoner at Lk 21:28.

The spoken body contains 1,171 words by TeXcount, excluding the title and this note. At 105--115 words per minute the estimated delivery is 10.2--11.2 minutes. The text was read through silently to check what a listener would need to remember from one paragraph to the next.

\section{References}

\begin{itemize}
\item \textbf{Appointed texts.} \work{Missale Romanum}, Vatican typical edition, 1962, First Sunday of Advent, pp. [1]--2, nos. 2--6 (Collect, Epistle, Gradual, Alleluia, Gospel) and 8--10 (Secret, Communion, Postcommunion). \url{https://media.churchmusicassociation.org/pdf/missale62.pdf}.
\item \textbf{Scripture.} Douay--Rheims Bible, Challoner revision. Lk 21:25--33, especially the words quoted from v. 28; Rom 13:8--14a, including the debt of love before the appointed lesson; Ps 24(25):3--4 and Ps 84:8, 13 (85:7, 12). \url{https://www.gutenberg.org/ebooks/1581}.
\item \textbf{Gregory the Great.} \work{Homiliae in Evangelia}, I, Homily 1, §§1--3, 5, Migne-derived Latin, on the disciples' expectation, the weary traveler who still wants his journey to continue and the storm Gregory's hearers remembered. The final assurance in the Gospel is discussed in §4.
\item \textbf{John Chrysostom.} \work{Homilies on Romans}, Homily 24, on Rom 13:11--14, particularly the runner near the prize and Paul's choice of “salvation.” Historical English in \work{Nicene and Post-Nicene Fathers}, first series, vol. XI.
\item \textbf{Augustine of Hippo.} \work{Enarrationes in Psalmos}, Psalm 84 (English heading Psalm LXXXV), §§10, 12, on Christ receiving from Mary the flesh in which he could die, and on the change in what a person delights to do. Historical English in \work{Nicene and Post-Nicene Fathers}, first series, vol. VIII. The homily applies this last distinction to the pleasure of injuring a neighbor.
\end{itemize}
